\section{Scope, semantics, and the main result}\label{sec:scope}

\subsection{What is being proved}

The two reference pages mix several logically different kinds of statement:
finite coefficient identities, convergence assertions, differential-equation
claims, integral representations, branch-dependent continuation formulas,
and global assertions about monodromy or modular functions.  Treating all of
these as though they had the same foundational content would conceal the main
issue.  In first-order \(\ipra\), variables range over natural numbers.  There
is no primitive sort of arbitrary real numbers, no quantification over all
complex functions, and no completeness axiom that manufactures limits without
rates.  We therefore use the following four-level interpretation.

\begin{definition}[Four proof levels]\label{def:levels}
A displayed classical statement is assigned one of the following readings.
\begin{enumerate}[label=\textsf{\Alph*.}]
\item \emph{Finite algebra.}  A statement about a finite product, a finite
sum, a coefficient, or a polynomial is translated literally into an
arithmetic identity or inequality between coded integers and rationals.
\item \emph{Effective analysis.}  A real or complex quantity is a fixed
primitive recursive approximation procedure with a fixed primitive recursive
Cauchy modulus.  An equality \(x=y\) means
\[
   \forall p\;\norm{x(p+s)-y(p+s)}\le 2^{-p}
\]
for an explicitly supplied precision shift \(s\).  Uniform statements carry
explicit compact-domain bounds and moduli.
\item \emph{Certified continuation.}  A branch or contour formula is read
relative to a finite code for a path, a branch of logarithm, positive distance
margins from singularities, and quantitative bounds on every local chart.
\item \emph{Global classical claim.}  A statement whose ordinary proof uses
unbounded path existence, compactness without data, an abstract Riemann
surface, or a global classification theorem is not silently reinterpreted as
an iPRA theorem.  We record the precise finite certificate that would suffice,
or mark the claim as outside the present package.
\end{enumerate}
\end{definition}

This distinction makes the strongest part of the theory unexpectedly cheap.
The hypergeometric function is controlled by one rational recurrence, and a
large fraction of its identities are consequences of that recurrence plus
uniqueness of a normalized formal solution of a second-order equation.

\subsection{Principal theorem}

Let a \emph{guarded rational parameter package} consist of rational codes
\(a,b,c,z\), together with proofs that every denominator used below is
nonzero.  A complex-rational package consists of pairs of rational codes and
uses the norm \(\norm{x+iy}=|x|+|y|\), whose comparisons are decidable.
For fixed packages, all operations in the following theorem are primitive
recursive.

\begin{theorem}[Constructive hypergeometric package]\label{thm:main-package}
The following families admit proofs in \(\ipra\), in the senses of
\cref{def:levels}.

\begin{enumerate}[label=(\roman*)]
\item[\StatusA--1.] The rising factorial, hypergeometric coefficient, finite
partial sum, generalized \({}_pF_q\) coefficient, and all their recurrence
relations are primitive recursive after totalized guarded division.

\item[\StatusA--2.] Symmetry in \(a,b\), termination for a nonpositive
integral numerator parameter, the explicit terminating polynomial,
coefficientwise differentiation of every finite order, the hypergeometric
ODE, and uniqueness of its normalized formal solution are provable by
quantifier-free induction.

\item[\StatusA--3.] The pullback--gauge calculations underlying Euler's
and Pfaff's transformations, the standard quadratic transformations, the six
Frobenius expressions, Kummer's 24 local forms, and all contiguous identities
on the source pages reduce to finite polynomial/rational certificates.
Equality of analytically continued branches is read at level \StatusB{} or
\StatusC{} with the displayed branch and path data.  Resonant logarithmic
second solutions are covered to every externally fixed truncation order; a
uniform order requires a jointly primitive recursive coefficient table.

\item[\StatusA--4.] Chu--Vandermonde, the terminating Jacobi and Legendre
representations, the generalized Chebyshev identity as a formal identity, the
binomial, logarithm, arcsine, incomplete-beta coefficient identity, and the
cubic-radical special case are provable by finite algebra and formal-series
uniqueness.

\item[\StatusB--1.] For each rational \(0\le r<1\), the series has a primitive
recursive geometric-tail modulus uniformly on \(\norm z\le r\), provided
quantitative lower bounds are supplied for denominator factors.  The sum is a
primitive recursively named complex number.  Termwise differentiation of any
fixed order, and of variable order under jointly primitive recursive data,
has an explicit modulus.

\item[\StatusB--2.] Euler's beta integral, its \([0,\infty)\) variant,
complete elliptic-integral specializations, and the incomplete-beta formula
hold for effective parameters with positive endpoint-exponent margins and a
branch/apartness certificate.  The proof is a finite quadrature argument plus
a geometric head--tail estimate.

\item[\StatusB--3.] Gauss's summation theorem in its Euler chamber, Kummer's
\(z=-1\) theorem, Gauss's second theorem and Bailey's \(z=1/2\) theorem follow
with explicit error schedules from Euler's integral, beta identities, and
quadratic transformations.  The terminating Vandermonde case has no analytic
hypotheses.

\item[\StatusB--4.] The coefficientwise confluent limit
\[
 {}_1F_1(a;c;z)=\lim_{h\to0}\F{a}{h^{-1}}{c}{hz}
\]
has an explicit compact-disk modulus after splitting into a finite head and a
uniform geometric tail.

\item[\StatusC--1.] The connection formula at \(z=1\), analytic continuation
along a rational polygonal path, Barnes's integral, and Gauss's continued
fraction hold relative to explicit nonresonance, branch, contour-separation,
tail, and denominator-apartness certificates described below.

\item[\StatusD--1.] The global Schwarz-triangle uniformization, the complete
finite-monodromy classification, the modular-lambda inversion statement, and
the modular proofs of all isolated MathWorld algebraic values are not claimed
as completed iPRA reductions here.  We prove the local algebraic differential
identities and give a finite certificate interface for the remaining global
input.
\end{enumerate}
\end{theorem}

\begin{proof}[Organization of the proof]
The finite core is developed in \cref{sec:arithmetic,sec:series,sec:ode}.
Effective convergence and differentiation are proved in
\cref{sec:convergence}.  Elementary and polynomial specializations are in
\cref{sec:special}; transformations and Kummer forms in
\cref{sec:transformations}; integrals and summations in
\cref{sec:integrals-sums}; contiguous relations and continued fractions in
\cref{sec:contiguous}; and continuation, Barnes contours, and the precise
boundary of the package in \cref{sec:continuation}.  Every analytic argument
is accompanied by a primitive recursive choice of truncation, mesh, or
precision.  The Lean compilation plan is in \cref{sec:formalization}.
\end{proof}

\subsection{Why theorem schemas are the right object}

A primitive recursive function may be represented in \(\ipra\) by a function
symbol together with defining equations.  It is important that we do not
quantify over codes for all primitive recursive programs and then assume a
primitive recursive universal evaluator: diagonalization rules that out.
Thus a theorem about a named approximation procedure \(A(n,p)\) is a theorem
schema, one first-order theorem for each fixed primitive recursive definition.
Numerical parameters of that definition may still vary uniformly.

This is exactly the format used by the surrounding ProveIt work.  Its ordinary
object theory has symbols for primitive recursive functions, their defining
equations, successor axioms, and quantifier-free induction.  Its checked
\code{Derivations} library packages theorems as sequents
\[
   ([]:\mathrm{List}(\mathrm{Fm}\;\mathrm{PTm}\;0))
      \;\vdash_{\mathrm{PRAAx}}\;\varphi,
\]
and transports them to the native scalar calculus through a bi-interpretation.
The present report supplies the mathematical derivations and the target
formulas; it does not claim that the full analytic development has already
been compiled into those sequents.

\subsection{Constructive content versus classical notation}

The notation \(\F{a}{b}{c}{z}\) is used at three levels:
\begin{enumerate}
\item the formal series \(\sum_{n\ge0}A_nz^n\), where equality means equality
of every coefficient;
\item the effective sum on a closed disk \(\norm z\le r<1\), represented by
partial sums and a tail modulus;
\item a continued branch obtained from a finite chain of local disks.
\end{enumerate}
The surrounding text always identifies the level.  A formal identity can be
stronger foundationally than an analytic one: it is valid over a localized
polynomial ring with no appeal to limits.  Conversely, a classical formula
that compares branches at distant points needs data that are invisible in the
printed equation.

\begin{remark}[No hidden excluded middle]
All comparisons used to choose a truncation or mesh are comparisons of coded
integers or rationals and hence decidable.  Bounded minimization over such a
test is primitive recursive.  We never infer a modulus from a merely
qualitative convergence assertion, and we never choose a zero or a path by an
unbounded existence principle.
\end{remark}

\subsection{A compact coverage map}

The source pages are covered in the following order.
\begin{center}
\small
\begin{tabularx}{\textwidth}{@{}p{.27\textwidth}p{.12\textwidth}X@{}}
\toprule
Source topic & Status & Constructive replacement \\
\midrule
Series, Pochhammer, termination & \StatusA & finite products, recurrence, finite polynomial \\
Differentiation & \StatusA/\StatusB & coefficient identity; analytic transfer by tail moduli \\
Elementary special cases & \StatusA/\StatusB & formal identities plus effective named functions \\
Hypergeometric ODE & \StatusA & coefficient recurrence and formal uniqueness \\
Six/24 local solutions & \StatusA/\StatusC & algebraic pullback/gauge certificates; branch data for analytic equality \\
Euler integral & \StatusB & rational quadrature with endpoint margins \\
Barnes integral & \StatusC & finite residue rectangle plus explicit vertical/horizontal tail bounds \\
Contiguous relations & \StatusA & coefficientwise rational identities \\
Continued fraction & \StatusA/\StatusC & exact finite convergents; convergence under apartness/contraction data \\
Euler--Pfaff/quadratic transforms & \StatusA/\StatusB & transformed ODE and normalized uniqueness \\
Gauss/Kummer/Bailey sums & \StatusA/\StatusB & finite Vandermonde or beta-integral evaluation \\
Schwarz maps and monodromy & \StatusC/\StatusD & local matrix/path certificates; no global classification claim \\
Modular-lambda and isolated modular values & \StatusD & exact displayed values audited numerically; modular proof not reconstructed \\
\bottomrule
\end{tabularx}
\end{center}
